\section{Explainability analysis}

Explainable artificial intelligence (XAI) is used here to investigate how the trained model exploits the waveform information that leads to the observed 
timing correction. While the coincidence timing resolution quantifies the performance of the model, it does not indicate which regions of the signal are 
responsible for that performance.

The purpose of the XAI analysis is therefore to identify where, in the aligned waveform, the model is most sensitive to the input and to relate these 
regions to the physical structure of the detector signal. This provides additional insight into the behavior of the model beyond the final CTR value 
and helps assess which parts of the waveform contain useful information for timing correction.

Understanding how the model uses the signal may also have practical implications beyond model interpretation. 
The identified informative waveform regions could suggest which signal features should be preserved or emphasized in future front-end electronics 
designs, while regions consistently assigned little importance may indicate where acquisition bandwidth or waveform duration could potentially be 
reduced. At the same time, the analysis can help clarify the intrinsic limits of the proposed correction approach by showing which information is 
actually available to the model in the measured waveform.

\subsection{Gradient-based temporal sensitivity}

The temporal regions used by the locally connected multilayer perceptron (LC-MLP) are examined with a gradient-based explainability analysis. For each event, the sensitivity of the predicted timing correction \(\widehat{y}\) to every normalized input sample is computed from the input gradient. The sample-wise importance is defined as the mean absolute gradient over the analyzed events and over the two detector branches:

\begin{equation}
    I_j
    =
    \frac{1}{2N}
    \sum_{n=1}^{N}
    \sum_{d=1}^{2}
    \left|
        \frac{\partial \widehat{y}_n}
             {\partial x_{n,d,j}}
    \right|,
\end{equation}

where \(x_{n,d,j}\) is sample \(j\) of detector \(d\) for event \(n\). Since the gradient is evaluated through the complete trained model, the sensitivity also includes the effect of the smooth bounded output transformation described previously.

The analysis is performed on up to 1024 events from the development population; the blind test population is not used to construct the XAI map. The resulting quantity should be interpreted as local model sensitivity rather than as a causal attribution: large values identify time samples for which small input variations produce a comparatively large change in the model output.

For visualization, the sample-wise importance is averaged within consecutive \SI{1}{ns} intervals. If \(B_k\) denotes the set of samples belonging to the \(k\)-th time interval,

\begin{equation}
    \overline{I}_k
    =
    \frac{1}{|B_k|}
    \sum_{j\in B_k} I_j .
\end{equation}

The interval values are then normalized to the largest value in the waveform, so that the displayed importance ranges from 0 to 1. In each plot, the upper panel shows one representative detector-waveform pair in physical voltage units, with the background shading indicating the normalized temporal sensitivity. The lower panel shows the corresponding \SI{1}{ns}-binned importance profile.

\subsection{Representative UC-board results}

Representative XAI maps for the UC board are shown for the energy-channel and timing-channel modes at \SI{48}{V} in Figs.~\ref{fig:xai-uc-energy} and~\ref{fig:xai-uc-timing}, respectively.

\reportfigure
    {../report/figures/xai/UC_energy_48V.pdf}
    {Gradient-based temporal sensitivity of the LC-MLP for the UC-board energy-channel study at \SI{48}{V}.}
    {fig:xai-uc-energy}

\reportfigure
    {../report/figures/xai/UC_timing_48V.pdf}
    {Gradient-based temporal sensitivity of the LC-MLP for the UC-board timing-channel study at \SI{48}{V}.}
    {fig:xai-uc-timing}

\subsection{Representative FBK-board results}

Representative XAI maps for the FBK board are shown for the energy-channel study at \SI{48}{V} and the timing-channel study at \SI{47}{V} in Figs.~\ref{fig:xai-fbk-energy} and~\ref{fig:xai-fbk-timing}, respectively.

\reportfigure
    {../report/figures/xai/FBK_energy_48V.pdf}
    {Gradient-based temporal sensitivity of the LC-MLP for the FBK-board energy-channel study at \SI{48}{V}.}
    {fig:xai-fbk-energy}

\reportfigure
    {../report/figures/xai/FBK_timing_47V.pdf}
    {Gradient-based temporal sensitivity of the LC-MLP for the FBK-board timing-channel study at \SI{47}{V}.}
    {fig:xai-fbk-timing}
